\section{Toward a Foundation for Mathematics}
\label{sec:foundations-roadmap}

The preceding sections give the scalar kernel.  They are not yet a full
foundation for mathematics.  A usable foundation must distinguish four kinds of
structure that ordinary prose often collapses:
\[
\text{value},\qquad
\text{semantic consequence},\qquad
\text{dependence},\qquad
\text{proof provenance}.
\]
Values may be Boolean, three-valued, or continuous.  Consequence is a global
relation over valuations or models.  Dependence is supplied by joints, joint
families, or copulas.  Proof provenance records which assumptions are actually
used in a derivation.  These structures answer different mathematical questions.

\paragraph{Semantic consequence.}
To prove ``$B$ from $A$'', the primitive object is not an implication truth
value.  It is the absence of a countermodel:
\[
  A\models_t B
  \quad\Longleftrightarrow\quad
  \forall v,\; t\le \val{A}_v \Rightarrow t\le \val{B}_v.
\]
The current labelled calculus captures this external reading.  Its formula
language has no internal conditional.  Its labels distinguish positivity,
certainty, and arbitrary thresholds.  This is enough to state graded consequence
cleanly, but not yet enough for a mature proof theory of mathematics.

\paragraph{Dependence.}
To update $A$ on evidence $B$, the primitive object is not the truth value of
$A$ or $B$ alone.  One needs a joint value or a family of possible joint values:
\[
  j=\text{value of }A\land B.
\]
The chain-rule judgment
\[
  c e=j
\]
then constrains the conditional value $c$.  If the joint is unknown, the result
is an interval or set of admissible conditionals.  This is the probability-side
lesson: fuzzy truth values alone are not enough; probability is logic enriched by
joint/dependence structure.

\paragraph{Proof provenance.}
To ask whether one theorem depends on another theorem, semantic consequence is
not enough.  If $B$ is already a theorem, then $A\models B$ for every $A$, but a
particular proof of $B$ may not use $A$.  A foundation for mathematics therefore
needs proof objects with dependency annotations.  The dependency question is:
which assumptions, lemmas, definitions, and side judgments occur in the proof
DAG?  This is a proof-theoretic question, not a truth-value question.

\paragraph{Hypothetical branches.}
Mathematical practice also uses inconsistent or counterfactual branches.  For
example, in the usual proof that $\sqrt2$ is irrational, one assumes the false
hypothesis that $\sqrt2$ is rational and derives a structured chain ending in
contradiction.  Classical explosion says that after contradiction everything
follows.  A useful foundation must preserve the dependency chain:
\[
  R \Rightarrow \text{coprime representation}
    \Rightarrow p^2=2q^2
    \Rightarrow p\text{ even}
    \Rightarrow q\text{ even}
    \Rightarrow \bot.
\]
The goal is not to make the false branch true.  The goal is to record what is
used, where inconsistency appears, and which conclusions are robust before
collapse.

\paragraph{Current status.}
The current Lean development already has the scalar conditioning kernel, labelled
external consequence, no-ex-falso examples, a first-order formula layer, equality
and quantifier interfaces, proof certificates, and a real-free checker path.  But
some proof-theoretic results are deliberately shallow: semantic consequence is
available as a proof rule, so completeness for that layer is a faithful interface
result rather than a deep syntactic completeness theorem.  This is honest and
useful, but it is not yet the final proof theory for mathematics.

\paragraph{Missing components.}
A serious foundation-building path should add the following modules or sections.

\begin{enumerate}[label=(F\arabic*),leftmargin=2.5em]
\item \textbf{Dependence contexts.}  A context should carry exact joints, interval
joints, or families of admissible joints.  Rules should propagate conditional
fibers through these contexts and mark conclusions as robust or dependence-sensitive.

\item \textbf{Proof provenance.}  Proof objects should expose the set or DAG of
used assumptions and side judgments.  This separates ``$A$ semantically entails
$B$'' from ``this proof of $B$ uses $A$.''

\item \textbf{Branch semantics.}  Hypothetical contexts such as $T+A$ should be
first-class, even when inconsistent.  The branch should support countermodel
elimination, local inconsistency detection, and no-explosion dependency tracing.

\item \textbf{Generative proof calculus.}  The labelled calculus should be extended
from a semantic-interface calculus to a more generative natural-deduction or sequent
calculus: introduction and elimination rules for the connectives, quantifiers,
equality, substitution, definitions, and induction where appropriate.  Semantic
imports should be isolated as oracle rules or removed from the core fragment.

\item \textbf{Robust collapse.}  Binary and three-valued collapse should happen only
after dependence-sensitive reasoning.  A collapsed statement should record whether
its status is robust for all admissible joints or depends on a chosen coupling.

\item \textbf{Mathematical library seed.}  The first target examples should be small:
parity, divisibility, order, finite sets, and simple contradiction proofs.  The
$\sqrt2$ irrationality proof is an ideal benchmark because it requires hypothetical
branching, equality/substitution, number theory lemmas, and provenance of the
contradiction.
\end{enumerate}

The resulting basis for mathematics is therefore not a single new truth table.
It is an architecture: formulas have values, theories generate model classes,
proofs carry provenance, dependence contexts carry joints or joint families, and
conditionals are constrained by the chain-rule fiber.  Classical mathematics is
recovered at certainty on crisp structures.  Probability enters when the joint
context is supplied or learned.  Ternary and binary views are then robust summary
layers, not places where dependence information is invented.
